% SPDX-License-Identifier: CC-BY-SA-4.0
% Author: Matthieu Perrin

\begingroup

\SetKwFunction{MPPatterns}{mp-patterns}
\SetKwFunction{RWPatterns}{rw-patterns}
\SetKwData{Reg}{reg}

\begin{exercice}[Mutual broadcast\footnote{[DMPR23] Mathilde Déprés, Achour Mostéfaoui, Matthieu Perrin, Michel Raynal. \emph{Send/Receive Patterns Versus Read/Write Patterns in Crash-Prone Asynchronous Distributed Systems.} DISC 2023}]

  On considère les deux algorithmes suivants, dans lesquels deux processus $p_0$ et $p_1$
  communiquent soit en s'envoyant des messages (algorithme \ref{algo:MPPatterns}), soit en
  lisant et écrivant dans des variables partagées (algorithme \ref{algo:RWPatterns}). 
  
  \begin{minipage}{.5\textwidth}
    \begin{algorithm}[H]
      \Method{$\MPPatterns(m)$}{
        \Broadcast{rb} $\textsc{msg}(i)$\;
        \Wait{$p_i$ has rb-delivered $\textsc{msg}(i)$} \;
        \Print $(p_i$ has rb-delivered $\textsc{msg}(1-i))$\;
      }
      \caption{Motifs en passage de messages}
      \label{algo:MPPatterns}
    \end{algorithm}
  \end{minipage}
  \begin{minipage}{.5\textwidth}
    \begin{algorithm}[H]
      $\Reg \leftarrow [\False, \False]$\;
      \Method{$\RWPatterns(m)$}{
        $\Reg[i] \leftarrow \True$\;
        \Print $(\Reg[1-i])$\;
      }
      \caption{Motifs en mémoire partagée}
      \label{algo:RWPatterns}
    \end{algorithm}
  \end{minipage}
  
  \begin{question}
  \item Décrivez toutes les sorties possibles des algorithmes \ref{algo:MPPatterns} et \ref{algo:RWPatterns}. 
  \end{question}

  On commence par définir la relation binaire $\lessdot$ sur les messages comme suit : pour toute paire de messages $m_i$ et $m_j$ diffusés
  par des processus $p_i$ et $p_j$ respectivement, on écrit $m_i \lessdot m_j$ si $p_j$ (l'émetteur de $m_j$) délivre $m_i$ avant $m_j$.
  De plus, pour tout message $m$, on note $\mathscr{P}_{\lessdot}(m) = \{m' | m' \lessdot m\}$ le passé de $m$ selon $\lessdot$. 

  
  Dans cet exercice, on s'intéresse à l'abstraction de communication \Broadcast{mutual},
  qui vérifie les mêmes propriétés que \Broadcast{rb}, en plus de la propriété d'ordre suivante : 
  \begin{description}
  \item[Mutual ordering :] $\lessdot$ est une relation totale. 
  \end{description}

  \begin{question}
    \item Reformulez, en français, la propriété Mutual ordering. 
  \item Démontrez la propriété suivante, liée au temps réel : si un processus $p_i$ \Broadcast{mutual} un message $m_i$,
    puis \Deliver{mutual} $m_i$ avant que $p_j$ \Broadcast{mutual} un message $m_j$, alors $p_j$ ne peut pas \Deliver{mutual} $m_j$ avant $m_i$. 
  \end{question}

  On considère maintenant l'algorithme~\ref{algo:MBroadcast} ci-dessous.
  
  \begin{algorithm}[H]
    \Method{$\Broadcast{mutual}(m)$}{
      \Broadcast{causal} $\textsc{init}(m)$\;
      \Wait{$p_i$ has causal-delivered $\textsc{ack}(m)$ \From at least $n-t-1$ processes} \;
      \Deliver{mutual} $m$ \From $p_i$\;
    }
    \Upon{$\textsc{init}(m)$ is causal-delivered \From $p_j$}{
      \If{$i\neq j$}{
        causal-\Send \textsc{ack}(m) \To $p_j$\;
        mutual-\Deliver $m$ \From $p_j$\;
      }
    }
    \caption{mutual-broadcast dans le modèle $\mathcal{CAMP}_{n, t}[causal-broadcast]$ (code pour $p_i$)}
    \label{algo:MBroadcast}
  \end{algorithm}

  \begin{question}
  \item Démontrez que l'algorithme~\ref{algo:MBroadcast} vérifie la propriété Mutual ordering dans le modèle \linebreak $\mathcal{CAMP}_{n, t}[causal-broadcast]$.
  \item On cherche maintenant à implémenter directement \Broadcast{mutual} dans le modèle  $\mathcal{CAMP}_{n, t}[\emptyset]$.
    Proposer une implémentation qui réduit la complexité au maximum. 
  \end{question}
\end{exercice}

\endgroup
\endinput
